\documentclass[a4paper, 12pt]{article}
\usepackage[utf8]{inputenc}
\usepackage[margin=1in]{geometry}
\usepackage{setspace}
\usepackage{xcolor}
\usepackage{hyperref}
\usepackage{lipsum}
\usepackage{fancyhdr}
\usepackage[para]{footmisc}

\renewcommand{\headrulewidth}{0pt}
\fancyfoot[R]{\thepage}
\singlespacing
\usepackage{titlesec}
\titleformat{\section}
  {\normalfont\large\bfseries}{\thesection}{1em}{}
\titleformat{\subsection}
  {\normalfont\normalsize\bfseries}{\thesubsection}{1em}{}
\titleformat{\subsubsection}
  {\normalfont\small\bfseries}{\thesubsubsection}{1em}{}
\titlespacing*{\section}{0pt}{1.5ex plus 1ex minus .2ex}{1ex plus .2ex}
\titlespacing*{\subsection}{0pt}{1.5ex plus 1ex minus .2ex}{0.8ex plus .2ex}
\titlespacing*{\subsubsection}{0pt}{1ex plus 1ex minus .2ex}{0.5ex plus .2ex}
\renewcommand{\thesubsection}{\arabic{subsection}}

\pagestyle{fancy}
\fancyhf{}
\fancyhead[L]{IK2221 VT26}
\fancyhead[R]{Lorenzo Deflorian (ldef@kth.se)}
%\fancyfoot[R]{\thepage}

\begin{document}

\section*{A High-Speed Load-Balancer Design with
Guaranteed Per-Connection-Consistency}
\subsection{Paper summary}
% Paper summary (for E level)
%% Problem being tackled and why it's hard
%% Key idea behind the approach
%% A bit more detail about the approach
%% The way the idea is evaluated, plus the key findings

This paper is about load balancing in large datacenters. The main issue is that you want to distribute connections evenly across servers, but at the same time you need all packets from the same connection to always go to the same server (PCC). In practice, existing solutions don't really manage to do both well. Hash-based load balancers, for example, either create noticeable imbalance or break connections when something changes, like when servers are added or removed.

The idea behind CHEETAH is actually pretty simple: instead of keeping track of connections inside the load balancer, it stores the mapping in a cookie that is attached to each packet. So when a packet comes in, the load balancer just reads the cookie and knows where to send it, without needing to store state. It kind of shifts the problem from the LB to the packets themselves.

They propose both a stateless and a stateful version. The stateless one can use basically any load balancing algorithm and still keep PCC, which is useful. The stateful version is more about efficiency, since it simplifies how mappings are stored compared to traditional approaches, even if its not completely new as a concept.

They implement the system both in software (FastClick) and on programmable hardware (Tofino), and test it with realistic workloads. The results show better load distribution and faster flow completion times (around 2-3x better than hash-based approaches), without adding much overhead, or at least not in a significant way.

\subsection{Contributions}
% Top 3 contributions of the paper  (extra, for C level)
%% Tie the contributions with impact.
%% The first paper to propose… | Makes it possible to… 
%% Presents a new technique for… | Dramatically improves performance of … 
%% Offers largely improved reliability… | Saves a large fraction of power for…

One important contribution is the idea of using cookies to guarantee PCC independently of the load balancing algorithm. This makes the system more flexible, since you're not tied to hashing anymore and you don't risk breaking connections when the system changes, which is a big deal in practice.

Another contribution is the design of both stateless and stateful versions. The stateless one removes the need for per-connection state, while the stateful one simplifies things compared to existing solutions. In particular, it avoids more complex data structures and ends up being faster, even if the improvement depends on the scenario.

Finally, they actually implement and test the system, which makes the paper more convincing overall. They show that using more advanced algorithms (like power-of-two-choices or weighted round robin) really improves performance when combined with CHEETAH, especially in terms of load balance and flow completion time, even if some assumptions are a bit ideal.

\subsection{Problems}
% Most glaring problem(s) with the work (up to 3) (extra, for A level)
%% Does the approach scale as the number of users increases?
%% Scaling: non-linear increase in resources | How does it work under a failure?
%% Does it miss some important cases? | Does it have too strong assumptions?
%% Does it have unrealistic deployment requirements?

There are a few things that felt a bit unclear or not fully explored.

First, deployment could be an issue. The whole approach depends on adding cookies to packets and having some additional components (like monitoring agents), but the paper mostly assumes programmable hardware or software load balancers. It's not obvious how easy it would be to use this in a more mixed or legacy environment, where not everything can be modified.

Second, the evaluation focuses more on normal operation and controlled changes, but doesn't go too deep into unexpected failures. For example, what happens if a server suddenly crashes or if different load balancers are temporarily inconsistent? These cases could affect PCC, but they're not really discussed in detail, or at least not enough.

Another point is about the need for frequent load updates when using more advanced policies. The paper shows good results, but it doesn't really quantify the cost of collecting all this information or how sensitive the system is to outdated data. At large scale, this could matter quite a bit and maybe reduce the benefits they show.

\end{document}